\chapter{Etapas da revisão bibliográfica}
\label{apend:2}

A revisão bibliográfica, por palavras-chave e \textit{snowballing}
\cite[p.~3-4]{wohlin2014}, teve duas rodadas, em setembro e outubro de
2026, no OpenAlex, no Crossref, na ACM Digital Library, na IEEE Xplore, na
SBC-OpenLib e na \gls{bdtd}.
Entram trabalhos revisados por pares e, da \gls{bdtd}, dissertações e
teses, em inglês ou português, de três tipos.
Os primeiros tratam de \gls{pr} ou da programação declarativa de interfaces
gráficas.
Os segundos comparam a compreensão e a usabilidade de uma ou de outra com
as dos \textit{callbacks} ou do \textit{Observer Pattern}.
Os terceiros aplicam as \glspl{dc} a linguagens e bibliotecas.
O \textit{snowballing} para a frente parte de dez trabalhos centrais e, na
segunda rodada, percorre todos os citantes de quatro deles em duas
iterações; uma terceira, parcial, vê as referências de um trabalho e os que
citam o 7GUIs, e não traz trabalho novo.
Para trás, percorre as referências de dois trabalhos.
A revisão não viu os citantes de \textcite{kiss2014}, que o OpenAlex não
indexa, e julgou só pelo título 280 registros sem resumo.

Entre os trabalhos incluídos, os mais próximos deste aplicam as \glspl{dc}
a linguagens e bibliotecas de interface ou de \gls{pr}:
\textcite{mernik2009} comparam o \gls{xaml} e o Windows Forms em C\#,
\textcite{kiss2014} compara JavaFX, ScalaFX, Scala.Rx, ReactFX e Elm, e
\textcite{zimmerle2025} avaliam o RxJS e o Bacon.js.
\textcite{grolaux2026} comparam, sem as \glspl{dc}, o \textit{async} e
\textit{await} com outras formas de tratar os eventos de interface.
Nenhum trabalho incluído avalia pelas \glspl{dc} a notação do React nem a
do Solid e do Angular com \textit{signals}, nem compara as três notações
nas mesmas tarefas; essa é a lacuna que a introdução declara.

As tabelas a seguir dão os números de cada etapa das duas rodadas.

\begin{table}[htbp]
\centering
\caption{\label{tab:revisaoRodada1}Primeira rodada (24 de setembro de 2026).}
\small
\begin{tabular}{p{0.55\textwidth}p{0.35\textwidth}}
\hline
Etapa & Registros \\
\hline
Buscas por palavras-chave (19 consultas) & 301 resultados \\
\textit{Snowballing} para a frente (10 trabalhos) & 359 citantes, de 932, com sobreposição \\
\textit{Snowballing} para trás (1 trabalho) & 82 referências \\
Títulos únicos triados & cerca de 530 \\
Lidos no resumo & cerca de 40 \\
Incluídos & 23 \\
\hline
\end{tabular}

Fonte: o autor.
\end{table}

\begin{table}[htbp]
\centering
\caption{\label{tab:revisaoRodada2}Segunda rodada (2 de outubro de 2026).}
\small
\begin{tabular}{p{0.55\textwidth}p{0.35\textwidth}}
\hline
Etapa & Registros \\
\hline
Citantes de 4 trabalhos, todos os anos & 2.359, com sobreposição \\
Citantes únicos desde 2013 & 1.230, dos quais 263 já vistos e 280 sem resumo \\
Iteração 1, incluídos por título e resumo & 89, dos quais 44 novos \\
Iteração 2: citantes de 13 trabalhos & 79, dos quais 70 inéditos; 1 incluído \\
Iteração 3: referências de 1 trabalho e trabalhos que citam o 7GUIs & 18; nenhum incluído \\
\gls{bdtd} (14 consultas) e SBC-OpenLib (10) & 508 títulos únicos; 2 incluídos, já na base \\
Buscas pela lacuna e por revisões de \gls{pr} no OpenAlex, no Crossref e na ACM DL & 1.528 títulos; nenhuma revisão nova \\
Candidatos novos & 12 \\
\hline
\end{tabular}

Fonte: o autor.
\end{table}
